\mainlanguage[de]

\setuplayout
  [backspace=2.0cm,
   height=middle,
   width=middle,
   leftmargin=1.8cm,
   footerdistance=0.0cm,
   rightmargindistance=0.2cm,
   rightmargin=1.8cm,
   bottom=0.5cm,
   leftmargindistance=0.2cm,
   header=2.1cm,
   headerdistance=0.5cm,
   footer=0.5cm,
   top=0cm,
   topspace=0.5cm,
   bottomspace=0.0cm]

\setupcolors[pagecolormodel=auto]

\setuphead
  [chapter]
  [number=no,
   style=,
   textstyle=\tfb\ss,
   before={\blank[small]},
   after={\blank[small]}]

\setuphead
  [section]
  [number=no,
   style=,
   textstyle=\tfa\ss,
   before={\blank[small]},
   after={\blank[small]}]

\definemeasure[headerimageheight][\dimexpr(\headerheight-2ex)]

\definelayer[sabel]

\startsetups[sabel]
\setlayerframed[sabel][x=0mm,y=0mm][frame=off,offset=none]
{%
   \setupTABLE[start][style=\tfx]
   \setupTABLE[r][each][height=0.85\baselineskip]
   \setupTABLE[column][1][align=right, width=.4\textwidth, loffset=-2pt, toffset=-5pt]%
   \setupTABLE[column][2][align={center,lohi}, width=.2\textwidth]%
   \setupTABLE[column][3][align=left, width=.4\textwidth, roffset=-2pt, toffset=-5pt]%
   \bTABLE[frame=off]
   \bTR \bTD Sabel-GBS           \eTD \bTD [nr=5] \dontleavehmode\externalfigure[../Logo/sabel.png][maxheight=\measure{headerimageheight}, frame=off]\dontleavehmode \eTD \bTD Schuljahr 2026/27                 \eTD \eTR
   \bTR \bTD Messtechnik         \eTD                                                                                                                                     \bTD Klasse M1A                        \eTD \eTR
   \bTR \bTD B.~Donner           \eTD                                                                                                                                     \bTD                                   \eTD \eTR
   \bTR \bTD \strut              \eTD                                                                                                                                     \bTD Name: \thinrules[n=1,width=2.7cm] \eTD \eTR
   \bTR \bTD Übungsblatt     \eTD                                                                                                                                     \bTD Einführung Messtechnik      \eTD \eTR
   \eTABLE}
\stopsetups

\setupbackgrounds[header][text][bottomframe=on,setups=sabel,background=sabel]
\setuppagenumbering[location={footer,middle},style=\bfc]

\setupfootertexts
  [Arbeitszeit: ca. 90 min][Messtechnik]
  [Messtechnik][Arbeitszeit: ca. 90 min]

\setdigitmode 3
\unexpanded\def\dn#1{\ifmmode\text{\digits{#1}}\else\digits{#1}\fi}

\setupmathematics[ampersand=text]
\setupformulas[align=middle]

\setupitemgroup[itemize][1][n,repeat][width=1.6em,distance=.4em]
\setupitemgroup[itemize][2][a,repeat][width=1.8em,distance=.4em]

\definemargindata
  [myinleft]
  [inleft]
  [align=flushleft]

\defineenumeration[exercise]
  [text=,
   alternative=myinleft,
   width=2em,
   distance=0.5em,
   style=\rm,
   headstyle=\rm,
   way=bytext,
   prefix=no,
   headcommand=\groupedcommand{}{}]

\defineframedtext
  [hinweis]
  [width=broad,
   frame=on,
   rulethickness=.6pt,
   corner=round,
   offset=.7em,
   background=color,
   backgroundcolor=lightgray,
   before={\blank[small]},
   after={\blank[medium]}]

\starttext
% Quelle: input/Übungen Einführung Messtechnik.pdf, Seiten 1–11.
% Reproduktion, keine Neuberechnung. Ungeklärte Scanstellen sichtbar markiert.
\define[1]\Loesung{{\it #1}\par\blank[small]}
\define[1]\ScanUnklar{{\tfxx [Scan unklar: #1]}}
\startMPinclusions
input "../include/mymacros.mp";
\stopMPinclusions
\startluacode
package.path = "../include/?.lua;" .. package.path
userdata = userdata or {}
userdata.plotmod = require("plotmod")
-- Aufgabe 11: relative geometrische Einheiten, keine numerischen Achsenwerte.
local T, A = 6, 2
local p = userdata.plotmod.create_plot()
p:set_window{ xmin=-0.5, xmax=1.7*T, ymin=-0.5, ymax=2*A+0.8,
              xunit="0.9cm", yunit="0.65cm" }
p:define_ticks{ x={{value=T,kind="major"}}, y={} }
p:define_labels{ x={}, y={} }
p:add_axes{ xlabel=[[\m{t}]], ylabel=[[\m{u}]],
            xlabel_placement="bottom_right", ylabel_placement="top",
            xaxis={label_distance="1mm",positive_end={arrow_gap="0mm"}},
            yaxis={label_distance="1mm",positive_end={arrow_gap="0mm"}} }
p:add_tex_label([[\m{T}]],T,0,{placement="bottom_right",distance="1mm"})
p:add_tex_label([[\m{\hat u}]],0,2*A,{placement="left",distance="1mm"})
p:add_function{ expression=string.format("%g*(1+sin(2*pi*x/%g))",A,T),
                samples=500,clip={xmin=0,xmax=1.5*T,ymin=0,ymax=2*A} }
-- Hilfslinie von Spannungsachse bis zum ersten Scheitelpunkt.
-- MP-Gerätekoordinaten einmal aus denselben unabhängigen Größen ableiten.
local xu = tex.sp("0.9cm")/65781.76
local yu = tex.sp("0.65cm")/65781.76
local ox,oy = .5*xu,.5*yu
p:add_raw_mp(string.format("draw (%g,%g)--(%g,%g) withpen pencircle scaled 0.4bp;",ox,oy+2*A*yu,ox+T/4*xu,oy+2*A*yu))
userdata.mischspannung = p:render()
\stopluacode

\chapter{Übungen – Einführung Messtechnik}
Übungsaufgaben
\section{I Analoge Messgeräte}
\startexercise
Es sollen Messungen mit einem analogen Spannungsmesser mit der
Genauigkeitsklasse \m{\mathrm{GK} = \dn{1.5}} im 20\unit{V} - Messbereich durchgeführt werden.
Für welche Spannungswerte ist eine relative Abweichung des Messergebnisses
vom tatsächlichem Wert von maximal \m{f = \pm 2\unit{Percent}} garantiert?
\blank[small]
\Loesung{Lösung: \m{U_{\mathrm{Mess}} >\! = 15\unit{V}}}
\stopexercise
\startexercise
Ein analoger Spannungsmesser mit der Genauigkeitsklasse \m{\mathrm{GK} = \dn{1.5}} liefert
beim Anlegen einer Spannung einen Messwert, der genau 40\unit{Percent} des
Vollausschlages beträgt.
\blank[small]
Bestimmen Sie den durch die Genauigkeitsklasse bedingten maximal möglichen
relativen Messfehler
\blank[small]
\Loesung{Lösung: \m{f = \pm \dn{3.75}\unit{Percent}}}
\stopexercise
\section{II Digitale Messgeräte}
\startexercise
Für ein 3 ½ stelliges DMM gibt der Hersteller die zusätzliche Angabe
\m{F = \pm (\dn{0.5}\unit{Percent}\ \text{v.MW} + 3\ \text{Digits})}.
\par
Wie groß darf ein Messwert im 20\unit{V} -Messbereich höchstens sein, damit ein
absoluter Fehler von \m{F = \pm 90\unit{mV}} nicht überschritten wird?
\blank[small]
\Loesung{Lösung: \m{\mathrm{MW} = 12\unit{V}}}
\stopexercise
\startexercise
Für ein 3 ½ stelliges DMM gibt der Hersteller die zusätzliche Angabe
\m{F = \pm (\dn{0.5}\unit{Percent}\ \text{v.MW} + 4\ \text{Digits})}.
\par
Wie groß darf ein Messwert im 200\unit{V} -Messbereich höchstens sein, damit ein
absoluter Fehler von \m{F = \pm \dn{1.15}\unit{V}} nicht überschritten wird?
\blank[small]
\Loesung{Lösung: \m{\mathrm{MW} = 150\unit{V}}}
\stopexercise
\startexercise
Ein digitaler Spannungsmesser hat eine Auflösung von 100\unit{microvolt}. Der Hersteller
gibt einen absoluten Fehler mit \m{\pm (\dn{0.25}\unit{Percent}\ \text{vom Messwert} + 1\ \text{Digits})} an. Der
Spannungsmesser zeigt einen Messwert von 300\unit{mV} an.
\blank[small]
Berechnen Sie, in welchem Bereich der angezeigte Messwert tatsächlich liegen
kann .
\blank[small]
\Loesung{Lösung: \m{\mathrm{MW} = 300\unit{mV} \pm \dn{0.85}\unit{mV}}}
\stopexercise
\page
\startexercise
Mit einem DMM wird im 2\unit{kiloohm}–Messbereich mit einer Auflösung von 1\unit{Ohm} und einer
Fehlergrenze von \m{\pm (\dn{1.6}\unit{Percent}\ \text{vom Messwert} + 2\ \text{Digits})} ein Wert von
\dn{0.152}\unit{kiloohm} angezeigt. Wird auf den 200\unit{Ohm} Messbereich mit einer Auflösung von
\dn{0.1}\unit{Ohm} umgeschaltet, der einen absoluten Fehler von
\m{\pm (\dn{1.6}\unit{Percent}\ \text{vom Messwert} + 2\ \text{Digits})} hat, so zeigt das Instrument einen Wert von
\dn{149.8}\unit{Ohm} an.
\blank[small]
Berechnen Sie für den jeweiligen Messbereich den relativen Fehler.
\blank[small]
\Loesung{Lösungen: \m{\mathrm{MB} = 2\unit{kiloohm} :\ f = \pm \dn{2.9}\unit{Percent}}\par
\hskip 5.4em \m{\mathrm{MB} = 200\unit{Ohm} :\ f = \pm \dn{1.73}\unit{Percent}}}
\stopexercise
\section{III Vergleich analoges und digitales Messgerät}
\startexercise
Ein analoges Voltmeter soll hinsichtlich der Genauigkeit mit einem DMM
verglichen werden. Der Vergleich soll für beide Geräte im 200\unit{V} Messbereich
durchgeführt werden.
\par
Das analoge Gerät besitzt eine Genauigkeitsklasse \m{\mathrm{GK} = \dn{0.5}}.
\par
Für das DMM gelten die Daten: \m{f_{\mathrm G} = \dn{0.5}\unit{Percent}\ ;\ f_{\mathrm Q} = 3\ ;\ Z = \dn{0.1}\unit{V}}
\blank[small]
\startitemize[a]
\sym{a)} Bei welchem Messwert arbeiten beide Instrumente gleich genau ?
\blank[small]
\Loesung{Lösung: \m{\mathrm{MW} = 140\unit{V}}}
\sym{b)} Welcher relativer Fehler ergibt sich bei diesem Messwert für beide Geräte ?
\blank[small]
\Loesung{Lösung: \m{f = \pm \dn{0.71}\unit{Percent}}}
\sym{c)} Welches von den beiden Geräten misst genauer wenn der unter a)
gefundene Wert überschritten wird ?
\blank[small]
\Loesung{Lösung: Das analoge Voltmeter}
\stopitemize
\stopexercise
\page
\section{IV Gleichricht- und Effektivwert}
\startexercise
Eine mit der Frequenz \m{f = 150\unit{Hz}} wiederkehrende Triggerpulsfolge, die mit
einem Drehspulmessgerät erfasst wird, ergibt einen Messwert von \m{U = 1\unit{V}}
\blank[small]
Wie groß ist die Impulsbreite, wenn die Impulsspannungshöhe \m{\hat u = 5\unit{V}} beträgt?
\blank[small]
\startlinecorrection
\midaligned{\startMPcode
numeric h, d, period;
h:=2.4cm; d:=1.1cm; period:=5.3cm;
pair O, D, P, E, H, Xend, Yend;
O:=(0,0); D:=O+(d,0); P:=O+(period,0); E:=P+(d,0);
H:=O+(0,h); Xend:=E+(1.7cm,0); Yend:=H+(0,.7cm);
pickup pencircle scaled .55bp;
drawarrow (O-(.5cm,0))--Xend;
drawarrow (O-(0,.5cm))--Yend;
draw H--(D+(0,h))--D--P--(P+(0,h))--(E+(0,h))--E
 withpen pencircle scaled .8bp;
label.top(textext("\m{\frac{U}{\unit{V}}}"),Yend);
label.lft(textext("\m{5}"),H);
label.bot(textext("\m{t_{\mathrm D}}"),D);
label.bot(textext("\m{T}"),P);
label.lrt(textext("\m{t}"),Xend);
\stopMPcode}
\stoplinecorrection
\blank[small]
\Loesung{Lösung : \m{t_{\mathrm D} = \dn{1.33}\unit{ms}}}
\stopexercise
\startexercise
Der folgende periodische Spannungsverlauf ist gegeben. Die Periodendauer sei
\m{T = 40\unit{ms}}.
\blank[small]
\startlinecorrection
\midaligned{\startMPcode
% Sprungpositionen nur geometrisch übernommen, ohne neue Zeitbeschriftung.
numeric period, high, low;
period:=6.3cm; high:=2.1cm; low:=-1.05cm;
pair O, Mid, Last, Jump, Peak, Xend, Yend, Highmark, Lowmark;
O:=(0,0); Mid:=O+(period/2,0); Last:=O+(period,0);
Jump:=.5[O,Mid]; Peak:=.5[Mid,Last]+(0,high);
Xend:=Last+(.9cm,0); Highmark:=O+(0,high); Lowmark:=O+(0,low);
Yend:=Highmark+(0,.6cm);
pickup pencircle scaled .55bp;
drawarrow (O-(1cm,0))--Xend;
drawarrow (Lowmark-(0,.4cm))--Yend;
draw ytick(1.5mm,Highmark);
draw xtick(1.5mm,Mid); draw xtick(1.5mm,Last);
draw (O-(2mm,2mm))--(O+(2mm,2mm));
draw (Last-(2mm,2mm))--(Last+(2mm,2mm));
draw Lowmark--(Jump+(0,low))--Jump--Peak--(xpart Peak,low)--(Last+(0,low))
 withpen pencircle scaled 1bp;
label.top(textext("\m{\frac{U}{\unit{V}}}"),Yend);
label.rt(textext("\m{20}"),Highmark+(2mm,0));
label.llft(textext("\m{-10}"),Lowmark);
label.bot(textext("\m{20}"),Mid-(0,2mm));
label.bot(textext("\m{40}"),Last-(0,2mm));
label.rt(textext("\m{\frac{t}{\unit{ms}}}"),Xend+(3mm,0));
\stopMPcode}
\stoplinecorrection
\blank[small]
Welchen Spannungswert zeigt ein Drehspulmessgerät mit vorgeschalteter
Gleichrichterschaltung an?
\blank[small]
\Loesung{Lösung: \m{U = 10\unit{V}}}
\stopexercise
\page
\startexercise
Eine Sägezahnspannung mit nachfolgender Austastzeit besitzt eine
Wiederholfrequenz von \m{f = 200\unit{Hz}} . Mit einem Drehspulmessgerät wird eine
Spannung von \m{U = \dn{0.5}\unit{V}} gemessen
\blank[small]
Wie groß ist die zeitliche Dauer \m{t_{\mathrm D}} des Sägezahns, wenn die maximale
Spannungshöhe \m{\hat u = 5\unit{V}} beträgt?
\blank[small]
\startlinecorrection
\midaligned{\startMPcode
numeric h, d, period;
h:=2.3cm; d:=1.7cm; period:=5.3cm;
pair O, D, P, E, H, Xend, Yend, Ga, Gb;
O:=(0,0); D:=O+(d,0); P:=O+(period,0); E:=P+(d,0);
H:=O+(0,h); Xend:=E+(1.2cm,0); Yend:=H+(0,.6cm);
Ga:=.35[D,P]; Gb:=.75[D,P];
pickup pencircle scaled .55bp;
draw (O-(.3cm,0))--Ga;
draw Ga--Gb dashed withdots;
drawarrow Gb--Xend;
drawarrow (O-(0,.5cm))--Yend;
draw O--(D+(0,h))--D;
draw P--(E+(0,h))--E;
draw ytick(1.5mm,H); draw xtick(1mm,P);
label.top(textext("\m{\frac{U}{\unit{V}}}"),Yend);
label.lft(textext("\m{5}"),H);
label.bot(textext("\m{t_{\mathrm D}}"),D);
label.bot(textext("\m{T}"),P);
label.lrt(textext("\m{t}"),Xend);
\stopMPcode}
\stoplinecorrection
\blank[small]
\Loesung{Lösung : \m{t_{\mathrm D} = 1\unit{ms}}}
\stopexercise
\startexercise
Eine sinusförmige Wechselspannung hat durch Überlagerung mit einer
Gleichspannung folgenden Verlauf:
\blank[small]
\startlinecorrection
\midaligned{\startMPcode
\ctxlua{context(userdata.mischspannung)}
\stopMPcode}
\stoplinecorrection
\blank[small]
Mit einem Drehspulmessgerät wurde diese Mischspannung gemessen und es ergab
sich ein Wert von u = \dn{7.5}\unit{V} .
\blank[small]
\startitemize[a]
\sym{a)} Welchen Spannungswert U hat die überlagerte Gleichspannung ?
\blank[small]
\Loesung{Lösung :U= \dn{7.5}\unit{V}}
\sym{b)} Bestimmen Sie bitte die Amplitude \m{\hat u} des dargestellten Spannungsverlaufs.
\blank[small]
\Loesung{Lösung : \m{\hat u = 15\unit{V}}}
\sym{c)} Mit einer weiteren Messung wird die reine Wechselspannung (ohne
Gleichspannungsanteil) wieder mit einem Drehspulmessgerät (mit
vorgeschaltetem Gleichrichter für beide Halbwellen) untersucht.
\par
Welcher Spannungswert wird jetzt angezeigt?
\blank[small]
\Loesung{Lösung: \m{U_{\mathrm{Gl}} = \dn{4.78}\unit{V}}}
\stopitemize
\stopexercise
\page
\startexercise
Bei einer Rechtecksignalspannung mit der Periodendauer \m{T = 80\unit{ms}} soll die
Impulshöhe \m{\hat u} sowie die Impulsbreite \m{t_{\mathrm D}} bestimmt werden.
\blank[small]
\startlinecorrection
\midaligned{\startMPcode
numeric h, d, period;
h:=2.4cm; d:=1.6cm; period:=4.8cm;
pair O, D, P, E, H, Xend, Yend, Ga, Gb;
O:=(0,0); D:=O+(d,0); P:=O+(period,0); E:=P+(d,0);
H:=O+(0,h); Xend:=E+(1.6cm,0); Yend:=H+(0,.8cm);
Ga:=.25[D,P]; Gb:=.75[D,P];
pickup pencircle scaled .55bp;
draw (O-(.8cm,0))--Ga;
draw Ga--Gb dashed withdots;
drawarrow Gb--Xend;
drawarrow (O-(0,.8cm))--Yend;
draw H--(D+(0,h))--D;
draw P--(P+(0,h))--(E+(0,h))--E;
label.rt(textext("\m{U}"),Yend);
label.lft(textext("\m{\hat u}"),H);
label.bot(textext("\m{t_{\mathrm D}}"),D);
label.bot(textext("\m{T}"),P);
label.lrt(textext("\m{t}"),Xend);
\stopMPcode}
\stoplinecorrection
\blank[small]
Zu diesem Zweck wird durch eine Messung der arithmetische Mittelwert und
durch eine weitere Messung der Effektivwert der Signalspannung bestimmt.
\blank[small]
\startitemize[a]
\sym{a)} Mit welchem analogen Messgerät wurde der arithmetische Mittelwert, und
mit welchem Instrument wurde der Effektivwert gemessen?
\blank[small]
\Loesung{Lösungen: Arithm.: Drehspul\par\hskip 5.4em Effekt.: Dreheisen}
\sym{b)} Bestimmen Sie die Impulshöhe \m{\hat u} und die Impulsbreite \m{t_{\mathrm D}}, wenn für den
arithmetischen Mittelwert u = \dn{1.25}\unit{V} und für den Effektivwert \m{U_{\mathrm{eff}} = \dn{2.5}\unit{V}}
gemessen wurden.
\blank[small]
\Loesung{Lösungen: \m{\hat u = 5\unit{V}\ ;\ t_{\mathrm D} = 20\unit{ms}}}
\stopitemize
\stopexercise
\section{V Widerstandsbestimmungen}
\startexercise
Ein unbekannter Widerstand \m{R_{\mathrm X}} soll durch eine Strom- und
Spannungsmessung ermittelt werden.
\par
Verwendet werden ein Spannungsmesser mit dem Innenwiderstand \m{R_{\mathrm{iU}} = 50\unit{kiloohm}}
\par
und ein Strommesser mit dem Innenwiderstand \m{R_{\mathrm{iI}} = \dn{0.5}\unit{Ohm}} wirksam durch die
jeweilig eingestellten Messbereiche.
\blank[small]
\startitemize[a]
\sym{a)} Zeichnen Sie die „Spannungsrichtige Schaltung“ mit Angabe der
Innenwiderstände der verwendeten Messgeräte.
\sym{b)} Durch die obige Schaltung wird eine Spannung von \m{U = \dn{9.97}\unit{V}} und ein
Strom \m{I = 20\unit{mA}} gemessen.
\par
Mit welchem Widerstandswert \m{\Delta R} muss der gemessene Wert korrigiert
werden, der sich durch den systematischen Fehler ergeben hatte?
\par
(Bitte möglichst genau rechnen.)
\blank[small]
\Loesung{Lösung: \m{\Delta R =}\ScanUnklar{Vorzeichen}\m{\dn{5.02}\unit{Ohm}}}
\stopitemize
\stopexercise
\page
\startexercise
Ein unbekannter Widerstand \m{R_{\mathrm X}} soll durch eine Strom- und Spannungsmessung
ermittelt werden.
\par
Verwendet werden ein Spannungsmesser mit dem Innenwiderstand \m{R_{\mathrm{iV}} = 80\unit{kiloohm}}
\par
und ein Strommesser mit dem Innenwiderstand \m{R_{\mathrm{iI}} = \dn{0.5}\unit{Ohm}}
\blank[small]
\startitemize[a]
\sym{a)} Zeichnen Sie die „Spannungsrichtige Schaltung“ mit den jeweiligen
Ersatzschaltbildern der verwendeten Messgeräte.
\sym{b)} Es wird eine Spannung von \m{U = 10\unit{V}} und ein Strom \m{I = \dn{1.55}\unit{mA}} gemessen.
\par
Wie groß ist der gemessene und der wahre Widerstandswert von \m{R_{\mathrm X}} ?
\par
(Bitte möglichst genau rechnen.)
\blank[small]
% Schwarzer Druck ohne k; rote handschriftliche k-Ergänzungen separat darüber.
\Loesung{Lösung:\m{R_{\mathrm{XM}} = \dn{6.452}\,\overset{\color[red]{k}}{\Omega},\ R_{\mathrm X} = \dn{7.018}\,\overset{\color[red]{k}}{\Omega}}}
\stopitemize
\stopexercise
\section{VI Messbereichserweiterung}
\startexercise
Bestimmen Sie die Widerstände \m{R_{\mathrm N}} und \m{R_{\mathrm V}} für eine Strom- und
Spannungsmessbereichserweiterung.
\blank[small]
Technische Daten des Messwerkes: \m{U_{\mathrm{max}} = 500\unit{mV}\ ;\ R_{\mathrm{iM}} = 500\unit{Ohm}}
\par
Gewünschte feste Messbereiche : \m{U = 10\unit{V}\ ;\ I = 20\unit{mA}}
\blank[small]
\Loesung{Lösungen: \m{R_{\mathrm N} = \dn{26.3}\unit{Ohm}\ ;\ R_{\text{[Index unklar]}} = 475\unit{Ohm}}}
\ScanUnklar{Index des zweiten R durch blaue Markierung überlagert.}
\stopexercise
\section{VII Relative Fehler}
\startexercise
Durch eine Strom- und Spannungsmessung soll sowohl eine
Widerstandsbestimmung als auch eine Leistungsbestimmung vorgenommen
werden.
\par
Das Strommessgerät arbeitet mit einer relativen Abweichung
\par
\m{f_{\mathrm I} = \Delta I / I = +\dn{0.004}}
\par
und das Spannungsmessgerät mit einer relativen Abweichung
\par
\m{f_{\mathrm U} = \Delta U / U = +\dn{0.006}}
\blank[small]
\startitemize[a]
\sym{a)} Welche relative Abweichung ergibt sich bei der Widerstandsbestimmung?
\blank[small]
\Loesung{Lösung: \m{f_{\mathrm R} = \dn{0.002}}}
\sym{b)} Bestimmen Sie nun die relative Abweichung bei der Leistungsbestimmung
durch die Strom-Spannungs-Beziehung (\m{P = U * I}),
\par
\hskip 2em die Strom-Widerstands-Beziehung (\m{P = R * I^2})
\par
und die Spannungs-Widerstands-Beziehung (\m{P = U^2 / R}).
\blank[small]
\Loesung{Lösung: \m{f_{\mathrm P} = \dn{0.01}}}
\stopitemize
\stopexercise
\page
\startexercise
Durch eine Strom- und Spannungsmessung soll sowohl eine
Widerstandsbestimmung als auch eine Leistungsbestimmung vorgenommen
werden.
\par
Das Strommessgerät arbeitet mit einer relativen Abweichung
\par
\m{f_{\mathrm I} = \Delta I / I = \dn{-0.004}}
\par
und das Spannungsmessgerät mit einer relativen Abweichung
\par
\m{f_{\mathrm U} = \Delta U / U = +\dn{0.005}}
\blank[small]
% Die Folge c), b) ist im Original so gedruckt.
\startitemize[a]
\sym{c)} Welche relative Abweichung ergibt sich bei der Widerstandsbestimmung?
\blank[small]
\Loesung{Lösung: \m{f_{\mathrm R} = \dn{0.9}\unit{Percent}}}
\sym{b)} Bestimmen Sie nun die relative Abweichung bei der Leistungsbestimmung
durch die Strom-Spannungs-Beziehung (\m{P = U * I}),
\par
\hskip 2em die Strom-Widerstands-Beziehung (\m{P = R * I^2})
\par
und die Spannungs-Widerstands-Beziehung (\m{P = U^2 / R}).
\blank[small]
\Loesung{Lösung: Immer \m{f_{\mathrm P} = \dn{0.1}\unit{Percent}}}
\stopitemize
\stopexercise
\startexercise
In dem abgebildeten Schaltungsauszug soll der Strom I bestimmt werden.
\blank[small]
\startlinecorrection
\midaligned{\startMPcode
pair Ctop, Cbot, Ntop, Nbot, N, Out, Aone, Atwo, Aout, Boxll, Boxur;
numeric r, sep;
r:=4mm; sep:=13mm;
Ctop:=(0,sep/2); Cbot:=Ctop-(0,sep);
N:=(22mm,0); Ntop:=(xpart N,ypart Ctop); Nbot:=(xpart N,ypart Cbot);
Out:=N+(21mm,0);
Aone:=.6[Ctop+(r,0),Ntop]; Atwo:=Aone-(0,sep); Aout:=.65[N,Out];
Boxll:=Nbot+(-1.5mm,-9mm); Boxur:=(xpart Out+5mm,ypart Ntop+9mm);
pickup pencircle scaled .55bp;
for C=Ctop,Cbot:
 draw fullcircle scaled (2r) shifted C;
 draw (C-(r+5mm,0))--(C-(r,0));
 draw (C-(r+11mm,0))--(C-(r+5mm,0)) dashed withdots;
endfor;
draw (Ctop+(r,0))--Ntop--Nbot--(Cbot+(r,0));
draw N--Out;
draw Out--(Out+(10mm,0)) dashed withdots;
drawarrow (Aone-(3mm,0))--Aone;
drawarrow (Atwo-(3mm,0))--Atwo;
drawarrow (Aout-(3mm,0))--Aout;
draw Boxll--(xpart Boxur,ypart Boxll)--Boxur--(xpart Boxll,ypart Boxur)--cycle dashed withdots;
label.top(textext("I1"),Aone+(0,2mm));
label.top(textext("I2"),Atwo+(0,2mm));
label.top(textext("I"),Aout+(0,2mm));
\stopMPcode}
\stoplinecorrection
\blank[small]
Der gestrichelt umrandete Teil ist jedoch nicht zugänglich. Aus diesem Grund
werden die Ströme I1 und I2 gemessen und anschließend der Strom I durch
Rechnung bestimmt.
\blank[small]
Es wurde gemessen: \m{\mathrm{I1} = 30\unit{mA} \pm 2\unit{Percent}\ ;\ \mathrm{I2} = \dn{-28.4}\unit{mA} \pm \dn{1.5}\unit{Percent}}
\blank[small]
\startitemize[a]
\sym{c)} Bestimmen Sie den Strom I .
\blank[small]
\Loesung{Lösung: \m{I = \dn{1.6}\unit{mA}}}
\sym{d)} Mit welchem relativen Fehler ist diese I -Bestimmung behaftet ?
\blank[small]
\Loesung{Lösung: \m{f = \pm \dn{64.1}\unit{Percent}}}
\stopitemize
\stopexercise
\page
\startexercise
Bei einer Quarzuhr soll ausschließlich die Periodendauer des Schwingquarzes
für die Ganggenauigkeit der Uhr ausschlaggebend sein. Bei der betrachteten
Uhr weicht die Quarzfrequenz vom Sollwert um \m{f = -\dn{0.001}\unit{Percent}} ab.
\blank[small]
Welche Zeitabweichung ergibt sich während einer Laufzeit von 14 Tagen
verglichen mit einer perfekt laufenden Uhr ?
\blank[small]
\Loesung{Lösung: \m{F = \dn{-12.096}\unit{s}} Die Uhr geht nach.}
\stopexercise
\section{VIII Normalverteilung}
\startexercise
Beim Herstellungsprozess von \dn{5.6}\unit{kiloohm}- Widerständen sind alle produzierten
Widerstände um den Sollwert R = \dn{5.6}\unit{kiloohm} mit einer Standartabweichung s = 110\unit{Ohm}
normalverteilt.
\blank[small]
\startitemize[a]
\sym{a)} Wieviel Prozent Ausschuss entstehen durch Aussortierung, wenn der
Hersteller eine Genauigkeit von \m{\Delta R/R = \pm 2\unit{Percent}} garantieren will.
\blank[small]
\Loesung{Lösung:}
\sym{b)} Skizzieren Sie die Verteilungsfunktion mit den Parametern x und s und
kennzeichnen Sie dort die Bereiche, die den Ausschuss darstellen.
\stopitemize
\ScanUnklar{Zusätzliche kleine Zeichen an x und s nicht sicher entzifferbar.}
\stopexercise

\stoptext
